\documentclass[11pt]{article}
\usepackage{estilo}
\addbibresource{referencias.bib}

\title{\LARGE\bfseries\color{tituloColor} Guia general para la redaccion de una tesis\\ en Ciencias de la Computacion}
\author{Plantilla conceptual --- no ligada a una nomenclatura fija de capitulos}
\date{}

\begin{document}
\maketitle

\begin{abstract}
\noindent
Este documento no es una tesis: es un \emph{molde argumentativo}, un ensayo que
explica que debe \textbf{existir} dentro de una tesis de Ciencias de la
Computacion para que esta sea correcta, autocontenida y defendible, con
independencia de los nombres que finalmente reciban sus capitulos. La tesis se
entiende aqui como lo describen \citet{booth2016craft}: una respuesta
justificada a una pregunta que le importa a una comunidad. Sobre esa base se
identifican cuatro funciones que todo trabajo de investigacion en computacion
debe satisfacer --- situar el problema, sostenerse a si mismo con la teoria
necesaria, distinguirse de lo ya hecho, y ofrecer una contribucion con
evidencia de que funciona --- y se cierra con la sintesis final. El objetivo
no es imponer una estructura rigida sino asegurar que, sin importar como se
titulen las secciones, la \emph{esencia} de cada funcion este presente.
\end{abstract}

\tableofcontents

% ============================================================
\section{Como leer y usar este molde}
% ============================================================

Una tesis de computacion no se organiza por convencion sino porque cada parte
cumple un rol argumentativo especifico. \citet{zobel2014writing} insiste en
que escribir en computacion es, ante todo, \emph{construir un argumento
verificable}: cada afirmacion debe apoyarse en evidencia, ya sea una
demostracion, un experimento o una referencia. Este documento traduce esa
idea en cuatro bloques funcionales que pueden llamarse de muchas maneras
(``Capitulo 1'', ``Introduccion'', ``Marco teorico'', ``Estado del arte'',
``Metodologia'', ``Resultados'', etc.) pero cuya \textbf{funcion} no es
negociable.

\begin{cajaNota}
Este documento esta escrito como ensayo argumentativo, no como tabla de
contenidos vacia. Cada bloque se explica indicando (i) que problema
resuelve dentro de la arquitectura global de la tesis, (ii) que elementos
minimos debe contener, (iii) por que la literatura metodologica exige esos
elementos, y (iv) cuando un elemento es opcional u obligatorio segun el tipo
de contribucion (algoritmo, modelo formal, sistema, estudio empirico).
\end{cajaNota}

La \cref{fig:pipeline} resume el flujo. Notese que las flechas no son solo de
avance: una tesis madura revisa constantemente sus primeros bloques a la luz
de lo que descubre en los ultimos, un proceso iterativo de refinamiento que
\citet{booth2016craft} describen como caracteristico de toda investigacion
seria, y que en la figura se representa mediante la flecha punteada de
retroalimentacion.

\begin{figure}[h]
\centering
\resizebox{\linewidth}{!}{%
\begin{tikzpicture}[
  node distance=1.1cm and 1.3cm,
  every node/.style={font=\small},
  bloque/.style={
    rectangle, rounded corners, draw=tituloColor, thick,
    minimum width=3.1cm, minimum height=1.5cm, align=center,
    text width=3.1cm
  }
]
  \node[bloque, fill=lightblue!70] (intro) {(a) Introduccion:\\ contexto, problema, objetivos};
  \node[bloque, fill=purple!25, right=of intro] (teoria) {(b) Fundamentos y estado del arte};
  \node[bloque, fill=orange!55, right=of teoria] (contrib) {(c) Contribucion y evidencia};
  \node[bloque, fill=green!35, right=of contrib] (concl) {(d) Conclusion};

  \draw[-{Latex[length=2.5mm]}, thick] (intro) -- (teoria);
  \draw[-{Latex[length=2.5mm]}, thick] (teoria) -- (contrib);
  \draw[-{Latex[length=2.5mm]}, thick] (contrib) -- (concl);

  \draw[-{Latex[length=2.5mm]}, thick, dashed, bordeAlerta]
    (contrib.south) .. controls +(0,-1.4) and +(0,-1.4) .. (intro.south)
    node[midway, below, yshift=-1.1cm, text width=6.5cm, align=center, font=\scriptsize\itshape]
    {iteracion: los resultados obligan a refinar el problema u objetivos};

  \begin{scope}[on background layer]
    \node[fit=(teoria)(contrib), fill=blue!8, rounded corners, inner sep=6pt,
          label={[font=\tiny, text=tituloColor]above:nucleo autocontenido de la contribucion}] {};
  \end{scope}
\end{tikzpicture}%
}
\caption{Flujo funcional de una tesis de Ciencias de la Computacion. Los
nombres de los capitulos pueden variar; las funciones (a)--(d) no.}
\label{fig:pipeline}
\end{figure}

% ============================================================
\section{Bloque (a): Introduccion --- situar el problema}
% ============================================================

La primera funcion de la tesis es \textbf{situar al lector}: por que este
problema importa, que se sabe y que no, y que se propone hacer al respecto.
\citet{creswell2018research} formaliza este bloque como la articulacion de
cuatro componentes que rara vez faltan en un proyecto de investigacion serio:
el planteamiento del problema, la revision preliminar de literatura que
motiva ese planteamiento, el proposito del estudio y las preguntas o
hipotesis que lo operacionalizan. Ninguno de estos componentes exige un
nombre de seccion especifico, pero los cuatro deben poder señalarse con el
dedo en algun lugar del documento.

\subsection{Contextualizacion y planteamiento del problema}

Debe explicarse el dominio de aplicacion o el area teorica, y dentro de ese
dominio, delimitar con precision el problema. \citet[cap.~2]{booth2016craft}
recomienda formular el problema como una tension no resuelta: algo que
\emph{deberia} poder hacerse (o entenderse, o demostrarse) y que actualmente
no se puede, no se sabe hacer bien, o no se sabe hacer de forma eficiente,
correcta o segura. Esta tension es lo que en la jerga de investigacion se
llama la \textbf{brecha} (\emph{gap}); situarla aqui evita que aparezca
recien en el capitulo de estado del arte, donde ya seria tardio.

\subsection{Objetivos, preguntas de investigacion e hipotesis}

Todo objetivo debe ser verificable por el propio diseño de la contribucion
que se presentara mas adelante: si el objetivo es ``diseñar un algoritmo mas
eficiente'', el bloque de evidencia debera demostrar esa eficiencia; si es
``demostrar una propiedad formal'', debera existir una prueba. Segun el tipo
de investigacion, esta seccion puede articularse como:

\begin{itemize}
  \item \textbf{Objetivo general y objetivos especificos}, tipico en trabajos
  de diseño o construccion de sistemas y algoritmos.
  \item \textbf{Preguntas de investigacion} (\emph{research questions}),
  frecuentes en estudios empiricos de ingenieria de software, siguiendo el
  patron de formulacion que describen \citet{wohlin2012experimentation}.
  \item \textbf{Hipotesis formales}, cuando el trabajo es de naturaleza
  cuantitativa o experimental y se someteran a prueba estadistica.
\end{itemize}

\begin{cajaAlerta}
Un objetivo que no puede fallar no es un objetivo cientifico. Si no existe
ninguna evidencia posible que pudiera mostrar que la contribucion \emph{no}
funciona, el objetivo esta mal planteado. Esta es una de las advertencias
mas repetidas por \citet{zobel2014writing}: la investigacion en computacion
debe ser \rojo{falsable o refutable en la practica}, no solo plausible en el
discurso.
\end{cajaAlerta}

\subsection{Alcance, supuestos y aporte esperado}

Cerrando el bloque, conviene anticipar brevemente en que consiste la
contribucion (sin desarrollarla aun) y cuales son sus limites deliberados:
que se asume, que queda fuera del alcance y por que. Esto no sustituye a la
conclusion; es una promesa que el resto de la tesis debe cumplir.

% ============================================================
\section{Bloque (b): Fundamentos necesarios y justificacion frente al estado del arte}
% ============================================================

Esta es la parte que garantiza que la tesis sea \textbf{autocontenida}: un
lector competente en computacion, pero no necesariamente experto en el
subtema, debe poder entender la contribucion sin salir del documento. Al
mismo tiempo, debe quedar claro por que la contribucion es necesaria frente
a lo que ya existe. Aunque en muchos programas estas dos tareas se separan en
capitulos distintos (``Marco teorico'' y ``Estado del arte''), conceptualmente
responden a una misma pregunta: \emph{?qué necesito saber para entender por
que mi contribucion es nueva y valida?}

\subsection{Fundamentos teoricos (autocontencion)}

Aqui se introduce el aparato conceptual minimo e imprescindible: definiciones
formales, notacion, modelos computacionales, estructuras de datos,
resultados previos que se usaran como caja negra, o el marco estadistico
necesario para interpretar los experimentos. \citet{denning1989computing}
argumentan que la disciplina de la computacion se sostiene sobre tres
paradigmas --- teoria, abstraccion (modelado) y diseño --- y que toda
contribucion puede describirse en terminos de al menos uno de ellos; esta
seccion debe dar al lector las herramientas del paradigma que corresponda a
la tesis.

\begin{cajaDefinicion}[: nivel de autocontencion exigible]
Una tesis es autocontenida cuando, tras leer este bloque, el lector puede
verificar por si mismo cada paso tecnico del bloque de contribucion sin
recurrir a fuentes externas para los conceptos \emph{centrales}. Puede
remitirse a fuentes externas para resultados \emph{auxiliares} bien
conocidos (por ejemplo, un teorema clasico de complejidad), siempre citando
su origen.
\end{cajaDefinicion}

Segun la naturaleza de la contribucion, el contenido minimo cambia:

\begin{itemize}
  \item \textbf{Contribucion algoritmica}: modelo de computo, definiciones de
  complejidad temporal/espacial, estructuras de datos empleadas.
  \item \textbf{Contribucion formal (teorema, sistema logico, semantica)}:
  definiciones, axiomas, lemas previos y convenciones notacionales completas,
  siguiendo las practicas de claridad que recomienda \citet{lamport1995write}
  para que cada paso de una demostracion sea verificable.
  \item \textbf{Contribucion de modelado} (arquitecturas, modelos
  probabilisticos, modelos de simulacion): supuestos del modelo, variables,
  y criterios de validez.
  \item \textbf{Contribucion de sistema/herramienta}: arquitectura general,
  decisiones de diseño y las tecnologias sobre las que se construye.
\end{itemize}

\subsection{Estado del arte y justificacion frente a lo existente}

No basta con listar trabajos relacionados: hay que \textbf{compararlos
criticamente} contra los objetivos planteados en el Bloque (a) para exponer
la brecha que la tesis llena. \citet{kitchenham2007guidelines} formalizan
este proceso mediante revisiones sistematicas de literatura, con criterios
explicitos de busqueda, inclusion y exclusion; aun cuando una tesis no
requiera una revision sistematica completa, adoptar su disciplina --- fuentes
trazables, criterios explicitos, cobertura representativa del area --- da
credibilidad a la afirmacion de novedad. \citet[cap.~2]{kuhn1962structure}
recuerda ademas que el valor de una contribucion cientifica se mide en
relacion con el paradigma vigente: una tesis no ataca al vacio, ataca las
limitaciones concretas de soluciones concretas.

\begin{cajaBuenaPractica}
Organice la comparacion en una tabla o taxonomia (por ejemplo, por enfoque,
supuestos, complejidad, o resultados reportados) y ubique explicitamente su
propia contribucion como una fila o celda adicional de esa misma tabla. Esto
hace visible la brecha en lugar de solo describirla en prosa.
\end{cajaBuenaPractica}

\begin{cajaAlerta}
Evite el ``estado del arte panoramico'': enumerar decenas de trabajos sin
conectarlos con los objetivos de la tesis. Cada trabajo citado en esta seccion
debe responder a una de dos preguntas: ?por que no resuelve el problema
planteado en el Bloque (a)? o ?que herramienta metodologica le presta a esta
tesis?
\end{cajaAlerta}

\subsection{La brecha como bisagra}

La brecha (\emph{gap}) identificada aqui debe coincidir exactamente con la
tension descrita en la Seccion~2.1. Si no coinciden, alguno de los dos
bloques esta mal planteado. Esta coherencia es, en la practica, el criterio
mas comun de rechazo en comites de tesis: objetivos que prometen algo que el
estado del arte revisado no muestra como necesario, o un estado del arte que
revela una brecha distinta de la que los objetivos atacan.

% ============================================================
\section{Bloque (c): Presentacion de la contribucion y evidencia}
% ============================================================

Este es el nucleo original de la tesis: aqui se presenta \emph{que se hizo} y
se demuestra \emph{que funciona}. La forma que tome depende enteramente del
tipo de contribucion, pero la exigencia de fondo es constante: toda
afirmacion de merito debe venir acompañada de evidencia de la clase adecuada
a esa afirmacion \citep{zobel2014writing}.

\subsection{Presentacion formal de la contribucion}

\begin{itemize}
  \item \textbf{Algoritmo}: especificacion precisa (pseudocodigo o
  formalismo equivalente), argumento de correctitud, y analisis de
  complejidad. Un algoritmo sin analisis de complejidad ni prueba de
  correctitud es una receta, no una contribucion cientifica.
  \item \textbf{Teorema o resultado formal}: enunciado preciso, demostracion
  completa y explicita. \citet{lamport1995write} propone estructurar toda
  demostracion como una jerarquia de pasos numerados, cada uno lo bastante
  pequeño para ser verificado de forma mecanica; esta estructura facilita
  enormemente la revision por el comite.
  \item \textbf{Modelo o arquitectura}: descripcion formal del modelo,
  justificacion de cada supuesto, y --- cuando corresponda --- un argumento
  de por que el modelo captura los aspectos relevantes del fenomeno.
  \item \textbf{Sistema o artefacto de software}: arquitectura, decisiones de
  diseño relevantes, y una descripcion suficiente para que el sistema sea
  reproducible por terceros.
\end{itemize}

\subsection{Diseño experimental (capitulo condicional)}

\begin{cajaNota}[: condicionalidad de la evaluacion experimental]
Si la contribucion hace afirmaciones de \emph{desempeño}, \emph{eficiencia
practica}, \emph{eficacia}, o cualquier propiedad no demostrable de forma
puramente analitica, entonces un capitulo o seccion de experimentos es
\textbf{obligatorio}. Si la contribucion es puramente formal y sus
propiedades quedan completamente establecidas por demostracion (por ejemplo,
un algoritmo cuya complejidad se deriva analiticamente y cuya correctitud se
prueba matematicamente), la evaluacion experimental es opcional y puede
limitarse a validar empiricamente el analisis teorico.
\end{cajaNota}

Cuando el capitulo experimental es necesario, su diseño no puede ser
improvisado. \citet{wohlin2012experimentation} y, antes, \citet[metodologia
seminal en ingenieria de software]{basili1986experimentation}, establecen
que un experimento defendible requiere, como minimo:

\begin{enumerate}
  \item \textbf{Definicion del objetivo experimental}, ligado directamente a
  las preguntas de investigacion del Bloque (a).
  \item \textbf{Variables independientes y dependientes}, y las metricas
  concretas con que se miden.
  \item \textbf{Diseño experimental} (grupos de control, lineas base de
  comparacion, aleatorizacion cuando aplique).
  \item \textbf{Amenazas a la validez} (interna, externa, de constructo y de
  conclusion), reportadas de forma explicita y no solo como formalidad.
  \item \textbf{Analisis estadistico} apropiado al tipo de dato recolectado.
\end{enumerate}

Para experimentos de rendimiento de sistemas (tiempos de ejecucion, uso de
memoria, escalabilidad), \citet{jain1991art} sigue siendo la referencia
estandar sobre metodologia de medicion, caracterizacion de carga de trabajo y
presentacion honesta de resultados: en particular, advierte contra reportar
un unico numero (por ejemplo, un promedio) sin exponer variabilidad,
condiciones del entorno de medicion ni supuestos de la carga usada.

\begin{cajaAlerta}
Una lista de baselines mal elegida invalida el capitulo entero. La linea
base debe ser el mejor competidor razonable disponible para el mismo
problema, no una version deliberadamente debil. Comparar solo contra una
implementacion \emph{naive} propia, cuando existen alternativas publicadas
del estado del arte, es una de las objeciones mas comunes en defensas de
tesis de corte experimental.
\end{cajaAlerta}

\subsection{Resultados e interpretacion}

Los resultados deben presentarse separados de su interpretacion: primero los
datos (tablas, graficas, medidas de dispersion), luego el analisis de lo que
significan en relacion con las hipotesis u objetivos. La \cref{tab:resumen}
resume, a modo de referencia cruzada, que evidencia corresponde a cada tipo
de contribucion.

\begin{table}[h]
\centering
\small
\rowcolors{2}{lightblue!25}{white}
\begin{tabular}{p{3.4cm}p{4.1cm}p{4.1cm}}
\toprule
\rowcolor{blue!30}
\textbf{Tipo de contribucion} & \textbf{Evidencia minima obligatoria} & \textbf{Evidencia adicional recomendada}\\
\midrule
Algoritmo & Prueba de correctitud + analisis de complejidad & Evaluacion empirica de rendimiento vs. baselines\\
Teorema / resultado formal & Demostracion completa & Ejemplos ilustrativos, casos limite\\
Modelo / arquitectura & Justificacion de supuestos + validacion cualitativa & Validacion cuantitativa contra datos reales\\
Sistema / herramienta & Descripcion reproducible de la arquitectura & Estudio de caso, benchmark, evaluacion de usuarios\\
Estudio empirico & Diseño experimental riguroso + analisis estadistico & Replicacion, analisis de amenazas a la validez\\
\bottomrule
\end{tabular}
\caption{Correspondencia entre tipo de contribucion y evidencia exigible.}
\label{tab:resumen}
\end{table}

% ============================================================
\section{Bloque (d): Conclusion}
% ============================================================

El cierre de la tesis no es un resumen mecanico: es el lugar donde se cierra
el argumento abierto en la Introduccion. \citet{phillips2010phd} señalan que
los comites evaluadores prestan especial atencion a que la conclusion
responda, de forma explicita, a cada objetivo o pregunta planteada al inicio;
dejar un objetivo sin respuesta clara --- incluso si la respuesta es
negativa o parcial --- es motivo frecuente de observaciones.

\subsection{Sintesis de la contribucion}

Se recapitula, en un parrafo, que se propuso, que se hizo y que se
demostro, evitando reintroducir informacion nueva. Debe leerse como el
cierre logico de la promesa hecha en la Seccion~2.3.

\subsection{Limitaciones}

Toda contribucion tiene fronteras. Declarar limitaciones no debilita la
tesis: la fortalece, porque demuestra que el autor entiende exactamente el
alcance de lo demostrado \citep{booth2016craft}. Las limitaciones deben
distinguirse de trabajos futuros: una limitacion es algo que el trabajo
\emph{actual} no cubre por diseño o por restriccion; el trabajo futuro es
una direccion que otros (o el mismo autor) podrian seguir despues.

\subsection{Trabajo futuro e implicancias}

Se identifican extensiones razonables, abiertas por los propios resultados o
por las limitaciones reconocidas. Cuando corresponda, tambien se discuten
implicancias practicas o teoricas mas amplias del hallazgo, conectando de
vuelta con la relevancia planteada en la Introduccion --- cerrando asi el
ciclo argumentativo completo de la tesis.

% ============================================================
\section{Consideraciones transversales de forma}
% ============================================================

Ademas del contenido, la forma en que se escribe cada bloque afecta su
credibilidad. \citet{zobel2014writing} y \citet{knuth1989mathematical}
coinciden en un conjunto reducido de principios que aplican a los cuatro
bloques por igual:

\begin{itemize}
  \item \textbf{Precision terminologica}: cada termino tecnico se define la
  primera vez que se usa y se usa de forma consistente en adelante.
  \item \textbf{Notacion unica}: un mismo simbolo no debe significar dos
  cosas distintas en capitulos distintos.
  \item \textbf{Verificabilidad}: cualquier afirmacion cuantitativa debe
  poder rastrearse hasta un dato, una demostracion o una cita.
  \item \textbf{Progresion logica}: cada seccion debe poder justificar, en
  una oracion, por que aparece despues de la anterior y antes de la
  siguiente.
\end{itemize}

\begin{cajaEjemplo}[: prueba de coherencia entre bloques]
Un ejercicio util antes de entregar la tesis es escribir, en una sola
oracion por bloque, la funcion que cumple: (a) ``expongo que X es un
problema no resuelto y me propongo lograr Y''; (b) ``explico lo necesario
para entender Y, y muestro que nadie lo ha logrado bajo las condiciones que
me propongo''; (c) ``presento Z y demuestro/mido que Z logra Y''; (d)
``confirmo que Z logro Y, dentro de los limites L, y sugiero W como
siguiente paso''. Si estas cuatro oraciones no encajan entre si, hay un
problema estructural que ninguna cantidad de redaccion cuidadosa en el
detalle podra ocultar ante un comite atento.
\end{cajaEjemplo}

% ============================================================
\section{Plantilla esqueleto de referencia}
% ============================================================

La siguiente estructura es un ejemplo de como los cuatro bloques
funcionales suelen materializarse en capitulos; los nombres son
intercambiables y pueden fusionarse o dividirse segun el programa de
posgrado o pregrado correspondiente.

\begin{cajaBuenaPractica}[: esqueleto orientativo]
\begin{enumerate}
  \item \textbf{Introduccion} --- contexto, problema, objetivos/hipotesis,
  alcance. \hfill\textit{(Bloque a)}
  \item \textbf{Fundamentos teoricos} --- definiciones, notacion, resultados
  previos usados como base. \hfill\textit{(Bloque b, parte teorica)}
  \item \textbf{Estado del arte} --- trabajos relacionados y brecha
  identificada. \hfill\textit{(Bloque b, parte comparativa)}
  \item \textbf{Contribucion propuesta} --- algoritmo, modelo, teorema o
  sistema, con su justificacion formal. \hfill\textit{(Bloque c, parte
  formal)}
  \item \textbf{Diseño experimental y resultados} --- \emph{solo si aplica
  segun la \cref{tab:resumen}}. \hfill\textit{(Bloque c, parte empirica)}
  \item \textbf{Conclusiones} --- sintesis, limitaciones, trabajo futuro.
  \hfill\textit{(Bloque d)}
\end{enumerate}
\end{cajaBuenaPractica}

Finalmente, conviene recordar la advertencia de \citet{wing2006computational}
sobre el pensamiento computacional: una tesis en esta disciplina no solo
resuelve un problema puntual, sino que aporta una forma reutilizable de
pensar sobre una clase de problemas. Mantener esa aspiracion presente en los
cuatro bloques --- y no solo en la introduccion --- es lo que distingue una
tesis solida de un reporte tecnico extenso.

\clearpage
\printbibliography

\end{document}
